% Lösungen Ebenen (zusammenbau v0.4, Heft)
\einheitenkopf{Ebenen}

\erg{37}{a) drei Punkte \quad b) $E\colon \vec x = (2 | 1 | 3) + r \cdot (1 | 0 | 4) + s \cdot (0 | 3 | 1)$, $r, s \in \mathbb{R}$ \quad c) $\vec x = (2 | 1 | -1) + r \cdot (2 | 1 | 1) + s \cdot (-2 | 2 | 2)$, $r, s \in \mathbb{R}$; $(-2 | 2 | 2)$ ist kein Vielfaches von $(2 | 1 | 1)$ (Faktor $-1$ und $2$)}
\erg{38}{a) $L\colon \vec x = (6 | 6 | 0) + r \cdot (0 | -3 | 3) + s \cdot (-3 | 0 | 3)$, $r, s \in \mathbb{R}$; aus I $s = 1$, aus II $r = 1$, III: $3 + 3 = 6$ ✓, $S$ liegt in $L$ \quad b) $O$, $A$ und $C$ liegen auf einer Geraden, weil $\overrightarrow{OC}$ ein Vielfaches von $\overrightarrow{OA}$ ist; $\overrightarrow{OB}$ ist kein Vielfaches von $\overrightarrow{OA}$, also liegt $B$ nicht auf dieser Geraden. Gerade und Punkt legen genau eine Ebene fest, und sie enthält alle vier Punkte. \quad c) $C$ liegt auf der Geraden durch $O$ und $A$, weil $\overrightarrow{OC}$ ein Vielfaches von $\overrightarrow{OA}$ ist; $\overrightarrow{OB}$ ist kein Vielfaches von $\overrightarrow{OA}$, also liegt $B$ nicht darauf. Die Gerade durch $O$ und $A$ und der Punkt $B$ legen genau eine Ebene fest, die alle vier Punkte enthält. \quad d) $L\colon \vec x = (1 | 0 | 3) + r \cdot (6 | 0 | 0) + s \cdot (0 | 5 | 3)$, $r, s \in \mathbb{R}$; aus I $r = 0{,}5$, aus II $s = 0{,}5$, III: $3 + 1{,}5 = 4{,}5$ ✓, $P$ liegt in $L$ \quad e) $L\colon \vec x = (0 | 0 | 3) + r \cdot (6 | 0 | 0) + s \cdot (0 | 4 | 2)$, $r, s \in \mathbb{R}$; aus I $r = 0{,}5$, aus II $s = 0{,}75$, III: $3 + 1{,}5 = 4{,}5 \ne 4$, $P$ liegt nicht in $L$}
\erg{39}{a) Koordinatenform \quad b) $\vec n = (-3 | -2 | 1)$ (jedes Vielfache ebenso) \quad c) $\frac{x}{6} + \frac{y}{3} + \frac{z}{2} = 1$, also $x + 2y + 3z = 6$}
\erg{40}{a) $d = 2 \cdot 3 + 3 \cdot 2 = 12$ ($y$ kommt nicht vor); Kontrolle mit $C$: $0 + 12 = 12$ ✓ \quad b) $\overrightarrow{CD} = (-3 | 3 | 0)$, $\overrightarrow{CS} = (-3 | -3 | 5)$, $\vec n = (5 | 5 | 6)$; $d = 5 + 5 + 30 = 40$ aus $S$, Probe mit $C$: $20 + 20 = 40$ ✓ \quad c) $\overrightarrow{BC} = (-1 | 3 | 0)$, $\overrightarrow{BS} = (-2 | 0 | 3)$, $\vec n = (3 | 1 | 2)$; $d = 9 + 1 + 6 = 16$ aus $B$, nicht aus dem Ursprung \quad d) $\overrightarrow{BC} = (0 | 8 | 0)$, $\overrightarrow{BE} = (-6 | 8 | 4)$, $\vec n = (2 | 0 | 3)$; $d = 12$ aus $B$: $L\colon 2x + 3z = 12$ ($y$ fehlt, $L$ ist parallel zur $y$-Achse) \quad e) $\vec n = (6 | 3 | 4)$ aus den Spannvektoren; $d = 24 + 3 = 27$: $E\colon 6x + 3y + 4z = 27$ \quad f) aus $F$: $c = 6$, aus $S$: $a = 3$, aus $M$: $9 + 3b + 12 = 18$, $b = -1$; $W\colon 3x - y + 6z = 18$ \quad g) $\vec x = (9 | 1 | 3) + r \cdot (-8 | 0 | 2) + s \cdot (-6 | 8 | 2)$, $r, s \in \mathbb{R}$; $\vec n = (4 | -1 | 16)$, $d = 36 - 1 + 48 = 83$ aus $W$ (einem Punkt der Plane): $E\colon 4x - y + 16z = 83$ \quad h) $\overrightarrow{AB} \times \overrightarrow{AC} = (80 | -4 | 0)$, gekürzt $\vec n = (20 | -1 | 0)$; $d = 20$ aus $A$: $E\colon 20x - y = 20$; $\vec x = (1 | 0 | 6) + r \cdot (2 | 40 | 0) + s \cdot (0{,}2 | 4 | 2)$, $r, s \in \mathbb{R}$}
\erg{41}{a) Parameterform \quad b) $(4 | -2 | 1)$ und z. B. $(-8 | 4 | -2)$ \quad c) $2x + y + 5z = 2 - 2 + 0$, also $E\colon 2x + y + 5z = 0$}
\erg{42}{a) $-3x + y + 4z = -6 + 5 + 4$, also $F\colon -3x + y + 4z = 3$ \quad b) $\overrightarrow{EF} = (-4 | 2 | 0)$, $\overrightarrow{ES} = (-4 | -1 | 2)$; $\vec n = (1 | 2 | 3)$; Probe: $-4 + 4 + 0 = 0$ und $-4 - 2 + 6 = 0$}
\erg{43}{a) $y$ \quad b) $y = 0$ \quad c) zur $y$-Achse: $y$ kommt nicht vor, $\vec n = (5 | 0 | 1)$ steht senkrecht auf der $y$-Richtung}
\erg{44}{a) $z$ fehlt und der Ursprung erfüllt die Gleichung: $E$ enthält die $z$-Achse, ist also nicht nur parallel zu ihr \quad b) $\overrightarrow{AB} = (4 | -2 | 0)$, $\overrightarrow{AD} = (-1 | 0{,}5 | 4)$, $\vec n = (1 | 2 | 0)$: die dritte Komponente ist 0, der Normalenvektor liegt waagerecht, die Wand steht senkrecht auf dem Boden}
\erg{45}{a) identisch \quad b) $F\colon 3x + y - 2z = 5$ \quad c) $F\colon 2x - y + 2z = 7$; $Q$: $6 - 1 + 2 = 7$ ✓, $Q$ liegt in $F$}
\erg{46}{a) $T\colon 3x + 2y + 3z = 21$ (mit $K$); $L$: $3 + 0 + 18 = 21$ ✓, $L$ liegt in $T$ (in $T$ einsetzen, nicht in $S$) \quad b) $(4 | 2 | 4) = 2 \cdot (2 | 1 | 2)$: parallel oder identisch; $(0 | 8 | 0)$ liegt in $L_1$, in $L_2$ ergibt sich $2 \cdot 8 = 16 \ne 14$ – die Ebenen sind echt parallel \quad c) $E_B\colon x - y + 6z = 6$; die Höhe $z = \frac{6 - x + y}{6}$ ist über $D$ am größten ($x = 0$, $y = 6$): $z = 2$ m (über $A$ nur 1 m) \quad d) Volumen $2 : 1$ heißt Höhen $2 : 1$, $M$ schneidet $BE$ nach einem Drittel von $B$ aus: $(1 | 1 | 0) + \frac{1}{3} \cdot (3 | 3 | 6) = (2 | 2 | 2)$; $M\colon x + 2y + 2z = 10$ (nicht die halbe Höhe) \quad e) Volumen $3 : 1$ heißt Höhen $3 : 1$, $M$ schneidet $AD$ nach einem Viertel von $A$ aus: $(1 | 1 | 0) + \frac{1}{4} \cdot (4 | 4 | 4) = (2 | 2 | 1)$; $M\colon 2x - y + 2z = 4$}
\erg{47}{a) parallel zu $E_1$: $E_2$ und $E_3$ ($E_3\colon x + 3z = 15$), nicht $E_4$; von unten nach oben: $E_1$, $E_2$, $E_3$}
\erg{48}{a) I und IV scheiden aus (Normalenvektor nicht kollinear zu $(2 | 0 | 5)$); für gleiches $x$ wächst $2x + 5z$ mit der Höhe, also mittlere Etage $2x + 5z = 50$, obere $2x + 5z = 60$; II läge unter der unteren Etage}
